\chapter{Manual de usuario}
\label{ap:manual}

\section{Configuración inicial del dispositivo}

La primera vez que se utiliza el dispositivo, es necesario realizar un proceso de configuración inicial que permite conectarlo a la red WiFi y vincularlo con la aplicación móvil.\\

\textbf{Paso 1:} Conectar el teléfono móvil a la red WiFi generada por el dispositivo (\textit{MyMeds-Setup}).\\

\textbf{Paso 2:} Abrir la aplicación móvil y acceder al apartado \textit{Configurar dispensador}.\\

\textbf{Paso 3:} Pulsar el botón \textit{Escanear QR}.\\

\textbf{Paso 4:} En la pantalla del dispositivo aparecerá un código QR que deberá ser escaneado desde la aplicación.\\%, como se muestra en la Figura~\ref{fig:manual_qr}.

% \begin{figure}[H]
%     \centering
%     \includegraphics[width=0.35\textwidth]{figs/manual_qr.jpg}
%     \caption{Escaneo del código QR durante la configuración inicial}
%     \label{fig:manual_qr}
% \end{figure}

\textbf{Paso 5:} Introducir las credenciales de la red WiFi (SSID y contraseña) donde se utilizará el dispositivo.\\

\textit{Nota:} En caso de disponer de redes WiFi de 2.4 GHz y 5 GHz, se debe seleccionar la red de 2.4 GHz, ya que el dispositivo no es compatible con redes de 5 GHz.\\

\textbf{Paso 6:} Confirmar el envío de las credenciales. Durante este proceso, tanto la aplicación como el dispositivo indicarán que se está realizando la sincronización.\\

\

\

\

\textbf{Paso 7:} Esperar a que finalice el proceso de sincronización.

Una vez completada la configuración:

\begin{itemize}
    \item La aplicación mostrará la pantalla de configuración de tomas.
    \item El dispositivo pasará a la pantalla principal, donde se muestra la hora y el día de la semana.
\end{itemize}

\section{Descripción del dispositivo}

El dispositivo es un pastillero inteligente diseñado para facilitar la gestión diaria de la medicación.

Incluye:

\begin{itemize}
\item Pantalla táctil para la interacción directa
\item Sistema automático de dispensación
\item Altavoz para avisos acústicos
\item Conectividad WiFi
\item Integración con aplicación móvil
\end{itemize}

\section{Uso del dispositivo}

\subsection{Pantalla principal}

Desde la pantalla principal puedes acceder a:

\begin{itemize}
\item Visualización de tomas del día
\item Historial de mediciones
\item Pulsómetro
\item Configuración de tomas
\end{itemize}

La navegación se realiza mediante la pantalla táctil.

\subsection{Dispensación de medicación}

El dispositivo dispensa automáticamente la medicación en los horarios configurados.

Cuando llega la hora y día de una toma:

\begin{itemize}
\item Se activa una señal sonora
\item Se muestra un aviso en pantalla
\item Se dispensa la medicación correspondiente
\end{itemize}

Cada compartimento del sistema corresponde a una toma concreta.

\subsection{Avisos acústicos}

El dispositivo incorpora un altavoz que emite un aviso cuando es necesario alertar al usuario.

El sonido ha sido diseñado para ser:

\begin{itemize}
\item Claramente audible
\item Suave y no molesto
\item Adecuado para uso diario
\end{itemize}

\section{Aplicación móvil}

La aplicación móvil permite configurar y gestionar el dispositivo de forma sencilla.

En este prototipo, la aplicación se proporciona ya instalada en el dispositivo móvil.

\subsection{Añadir una toma}

Para crear una nueva toma:

\begin{itemize}
\item Pulsa en “Añadir toma”
\item Selecciona la hora
\item Elige los días de repetición
\item Introduce la medicación
\item Activa o desactiva el recordatorio
\item En caso de quererlo, elige cuantos minutos antes quieres un aviso previo a la toma
\item Guarda la configuración
\end{itemize}

\subsection{Editar una toma}

Puedes modificar cualquier toma existente pulsando ``editar'' en la toma:

\begin{itemize}
\item Cambiar la hora
\item Modificar los días
\item Activar o desactivar el recordatorio
\item Modificar el aviso previo
\end{itemize}

\subsection{Eliminar una toma}

Para eliminar una toma, selecciona la opción ``Eliminar'' de la toma. Aparecerá un mensaje de confirmación para proceder con la eliminación.

\subsection{Sincronización con el dispositivo}

Cada vez que se realiza un cambio en la aplicación:

\begin{itemize}
\item Se envía la configuración completa al dispositivo
\item Se actualizan todas las tomas almacenadas
\end{itemize}

El sistema no realiza actualizaciones parciales, sino que reemplaza la configuración anterior por la nueva.

\section{Funcionamiento del sistema}

El dispositivo almacena las tomas en su memoria interna y funciona de forma autónoma.\\

El sistema comprueba continuamente la hora actual y determina si existe una toma programada.\\

En caso afirmativo:

\begin{itemize}
\item Se activa el aviso acústico
\item Se muestra la información en pantalla
\item Se realiza la dispensación
\end{itemize}

\section{Funcionamiento sin conexión}

Una vez configurado, el dispositivo no requiere conexión continua con la aplicación.

\begin{itemize}
\item Mantiene las tomas almacenadas
\item Sigue funcionando con normalidad
\item No recibe cambios hasta una nueva sincronización
\end{itemize}

Si el dispositivo se queda sin corriente, no será necesario volver a sincronizarlo con la aplicación.\\

Solo será necesario hacerlo si se borra el caché o datos de la aplicación o si se borra la memoria del dispositivo.

\section{Sensor biomédico}

El dispositivo incluye un sensor para la medición de parámetros como la frecuencia cardíaca y la saturación de oxígeno.\\

Para realizar una medición:

\begin{itemize}
\item Coloca el dedo sobre el sensor
\item Mantén el dedo inmóvil durante unos segundos
\item Espera a que aparezca el resultado en pantalla
\end{itemize}

\section{Recomendaciones de uso}

Para un correcto funcionamiento:

\begin{itemize}
\item Mantén el dispositivo en un entorno limpio y seco
\item Evita golpes o caídas
\item Utiliza una fuente de alimentación adecuada 5V/3A
\item Revisa periódicamente el estado del dispensador
\end{itemize}

\section{Problemas frecuentes}

\textbf{El dispositivo no enciende}
Comprueba la conexión del cable USB-C.

\textbf{No aparecen las tomas}
Realiza una modificación en la aplicación para forzar la sincronización.

\textbf{El dispositivo no tiene configuración}
Vuelve a sincronizar desde la aplicación móvil.

\textbf{No se escucha el aviso}
Verifica el funcionamiento del altavoz.

\textbf{No conecta al WiFi}
Reinicia el dispositivo e introduce de nuevo los datos de la red.

\section{Información adicional}

El sistema ha sido diseñado para ofrecer una experiencia sencilla, intuitiva y accesible, especialmente orientada a usuarios que requieren apoyo en la gestión de su medicación diaria.\\

La combinación del dispositivo y la aplicación permite una gestión completa, automatizada y fiable de las tomas.
